\documentclass[11pt]{article}
\usepackage[margin=1in]{geometry}
\usepackage{amsmath,amssymb}
\setlength{\parindent}{0pt}
\setlength{\parskip}{6pt}

\title{STT 465: Bayesian Data Analysis\\Homework 4}
\author{Mihir Patel}
\date{}

\begin{document}
\maketitle

%==================================================
\section*{Question 1}

$X_1,\dots,X_n \mid \sigma^2 \sim N(\mu,\sigma^2)$, \quad $\mu = 5$, \quad $\lambda = \frac{1}{\sigma^2}$, \quad $SS_\mu = \sum (x_i-\mu)^2$, \quad $\lambda \sim \text{Gamma}(a_0,b_0)$

\subsection*{1(a)}
\begin{align*}
f(x\mid\lambda) &= \prod_{i=1}^{n} \left(\frac{\lambda}{2\pi}\right)^{1/2} \exp\left[-\frac{\lambda}{2}(x_i-\mu)^2\right] \\
&= (2\pi)^{-n/2}\,\lambda^{n/2} \exp\left[-\frac{\lambda}{2}\sum (x_i-\mu)^2\right]
\end{align*}
Dropping the constant $(2\pi)^{-n/2}$:
\[
\boxed{f(x\mid\lambda) \propto \lambda^{n/2}\exp\left(-\frac{\lambda\, SS_\mu}{2}\right)}
\]

\subsection*{1(b)}
\begin{align*}
\pi(\lambda\mid x) &\propto f(x\mid\lambda)\,\pi(\lambda) \propto \lambda^{n/2} e^{-\lambda SS_\mu/2}\cdot \lambda^{a_0-1}e^{-b_0\lambda} \\
&= \lambda^{(a_0+n/2)-1}\exp\left[-\left(b_0+\frac{SS_\mu}{2}\right)\lambda\right]
\end{align*}
\[
\boxed{a_n = a_0 + \frac{n}{2}, \qquad b_n = b_0 + \frac{SS_\mu}{2}}
\]
These are the two parameters since this is the kernel of $\lambda\mid x \sim \text{Gamma}(a_n,b_n)$.

\subsection*{1(c)}
$a_0 = 2$, $b_0 = 3$, $n = 8$, $SS_\mu = 18$
\[
a_n = 2 + \frac{8}{2} = 6, \qquad b_n = 3 + \frac{18}{2} = 12
\]
\[
\lambda\mid x \sim \text{Gamma}(6,12), \qquad E(\lambda\mid x) = \frac{a_n}{b_n} = \frac{6}{12} = 0.5
\]

\subsection*{1(d)}
Jeffreys prior: $\pi(\lambda)\propto \lambda^{-1}$, i.e.\ $a_0, b_0$ approaching 0.
\[
\pi(\lambda\mid x) \propto \lambda^{n/2-1} e^{-\lambda SS_\mu/2}
\;\Rightarrow\;
\lambda\mid x \sim \text{Gamma}\left(\frac{n}{2},\frac{SS_\mu}{2}\right) = \text{Gamma}(4,9)
\]
\[
E(\lambda\mid x) = \frac{4}{9} \approx 0.444
\]
The Jeffreys posterior mean is the same as the MLE $\hat\lambda = \frac{n}{SS_\mu} = \frac{8}{18} = 0.444$. It is driven entirely by the data. The Gamma(2,3) prior has mean $\frac{2}{3} = 0.667$ and acts like $2a_0 = 4$ extra observations, so it ends up pulling the posterior mean up from 0.444 to 0.5, between the prior mean and the MLE.

\subsection*{1(e)}
The posterior prediction for a future $\tilde{x}$ averages the normal over the posterior for $\sigma^2$:
\[
p(\tilde{x}\mid x) = \int_0^\infty N\left(\tilde{x}\mid \mu, \tfrac{1}{\lambda}\right)\,\text{Gamma}(\lambda\mid a_n,b_n)\,d\lambda
\]
This is a scale mixture of normals, giving the Student $t$ distribution with $2a_n$ degrees of freedom. It has heavier tails than any single normal.
\[
\text{Var}(\tilde{x}\mid x) = E(\sigma^2\mid x) = \frac{b_n}{a_n-1}
\]
Plugging in one fixed $\sigma^2$ ignores our uncertainty for the variance. Averaging over the plausible values of $\sigma^2$, including the large ones, adds uncertainty and spreads our predictive out. Here the predictive is $t_{12}(5, 2)$ with variance $\frac{12}{5} = 2.4$.

%==================================================
\section*{Question 2}

$X\mid\mu,\sigma^2 \sim N(\mu,\sigma^2)$, \quad $\mu\mid\sigma^2 \sim N\left(\mu_0,\frac{\sigma^2}{p_0}\right)$, \quad $\sigma^2 \sim IG\left(\frac{\nu_0}{2},\frac{SS_0}{2}\right)$, \quad $SS = \sum(x_i-\bar{x})^2$

\subsection*{2(a)}
\[
\mu\mid\sigma^2,x \sim N\left(\mu_n,\frac{\sigma^2}{p_n}\right), \qquad \sigma^2\mid x \sim IG\left(\frac{\nu_n}{2},\frac{SS_n}{2}\right)
\]
\[
p_n = p_0 + n, \qquad \mu_n = \frac{p_0\mu_0 + n\bar{x}}{p_0+n}, \qquad \nu_n = \nu_0 + n, \qquad SS_n = SS_0 + SS + \frac{p_0 n}{p_0+n}(\bar{x}-\mu_0)^2
\]

\subsection*{2(b)}
$p_0 = 2$, $\mu_0 = 4$, $\nu_0 = 4$, $SS_0 = 12$, $n = 6$, $\bar{x} = 5$, $SS = 15$
\begin{align*}
p_n &= 2 + 6 = 8 \\
\mu_n &= \frac{2(4) + 6(5)}{8} = \frac{38}{8} = 4.75 \\
\nu_n &= 4 + 6 = 10 \\
SS_n &= 12 + 15 + \frac{2\cdot 6}{8}(5-4)^2 = 27 + 1.5 = 28.5
\end{align*}

\subsection*{2(c)}
\[
\mu\mid\sigma^2,x \sim N\left(4.75,\frac{\sigma^2}{8}\right)
\]
\[
\sigma^2\mid x \sim IG\left(\frac{10}{2},\frac{28.5}{2}\right) = IG(5,\,14.25)
\]

\subsection*{2(d)}
\[
\mu_n = \frac{p_0}{p_0+n}\mu_0 + \frac{n}{p_0+n}\bar{x} = \frac{2}{8}(4) + \frac{6}{8}(5) = 1 + 3.75 = 4.75
\]
The weight on $\bar{x}$ is $\frac{n}{p_0+n}$, which increases to 1 as $n$ approaches infinity. The prior weight is $\frac{p_0}{p_0+n}$. With more data, the posterior mean is driven almost entirely by our sample mean. The prior acts like $p_0$ pseudo-observations centered at $\mu_0$.

\subsection*{2(e)}
$SS_0$ is our prior sum of squares and $SS$ measures spread within our data. $\frac{p_0 n}{p_0+n}(\bar{x}-\mu_0)^2$ measures how far the sample mean is from the prior mean. If those two conflict, $SS_n$ inflates, so our posterior $\sigma^2$ shifts toward larger values, allowing more uncertainty. $\frac{p_0 n}{p_0+n}$ is the combination of the two sample sizes.

%==================================================
\section*{Question 3}

$x = (20,12,8)$, \quad $n = 40$, \quad $x\mid\omega \sim \text{Mult}(40,\omega)$, \quad $\omega \sim \text{Dir}(2,2,2)$

\subsection*{3(a)}
\[
f(x\mid\omega) \propto \omega_1^{20}\,\omega_2^{12}\,\omega_3^{8}
\]
\[
\pi(\omega) \propto \omega_1^{2-1}\,\omega_2^{2-1}\,\omega_3^{2-1} = \omega_1\,\omega_2\,\omega_3
\]
\[
\pi(\omega\mid x) \propto \omega_1^{21}\,\omega_2^{13}\,\omega_3^{9} = \omega_1^{22-1}\,\omega_2^{14-1}\,\omega_3^{10-1}
\]
\[
\boxed{\omega\mid x \sim \text{Dir}(2+20,\,2+12,\,2+8) = \text{Dir}(22,14,10)}
\]

\subsection*{3(b)}
\begin{align*}
E(\omega_1\mid x) &= \frac{22}{46} \approx 0.4783 \quad \text{(car)} \\
E(\omega_2\mid x) &= \frac{14}{46} \approx 0.3043 \quad \text{(bus)} \\
E(\omega_3\mid x) &= \frac{10}{46} \approx 0.2174 \quad \text{(bike)}
\end{align*}

\subsection*{3(c)}
\[
P(\tilde{y} = \text{bike}\mid x) = \int \omega_3\,\pi(\omega\mid x)\,d\omega = E(\omega_3\mid x) = \frac{10}{46} \approx 0.2174
\]

\subsection*{3(d)}
Posterior: $\text{Dir}(21,13,9)$
\[
E(\omega_1\mid x) = \frac{21}{43} \approx 0.4884
\]
This is larger than 0.4783 and closer to our sample proportion $\frac{20}{40}$. Dir(2,2,2) is worth 6 pseudo-counts spread evenly, so it shrinks every proportion toward $\frac{1}{3}$.

\subsection*{3(e)}
Use $\alpha = (22,14,10)$. Since $Y_k \sim \text{Gamma}(\alpha_k,1)$ and $S = \sum Y_k$,
\[
\left(\frac{Y_1}{S},\frac{Y_2}{S},\frac{Y_3}{S}\right) \sim \text{Dir}(\alpha_1,\alpha_2,\alpha_3)
\]

% ---- PASTE R CODE FOR 3(e) BELOW ----

\begin{verbatim}
# STT 465 HW4 - Question 3(e): Dirichlet posterior via Gamma normalization
set.seed(465)
U <- 10000
alpha <- c(22, 14, 10)   # posterior Dir(2+20, 2+12, 2+8)

# Draw independent Gamma(alpha_k, 1) variables, then normalize
y1 <- rgamma(U, shape = alpha[1], rate = 1)
y2 <- rgamma(U, shape = alpha[2], rate = 1)
y3 <- rgamma(U, shape = alpha[3], rate = 1)
S  <- y1 + y2 + y3
w1 <- y1 / S; w2 <- y2 / S; w3 <- y3 / S

c(mean(w1), mean(w2), mean(w3))       # check vs. 22/46, 14/46, 10/46
mean(w1 > w2)                         # P(w1 > w2 | x)
quantile(w1 - w2, c(0.025, 0.975))    # 95% equal-tailed CI for w1 - w2
\end{verbatim}
Output:
\begin{verbatim}
[1] 0.4790278 0.3045330 0.2164392
[1] 0.9131
       2.5%       97.5%
-0.07420041  0.41572194
\end{verbatim}


$P(\omega_1 > \omega_2\mid x) \approx 0.913$. The 95\% equal-tailed credible interval for $\omega_1-\omega_2$ is $(-0.074,\ 0.416)$, which means car is very likely more popular than bus, but our interval still includes 0.

%==================================================
\section*{Question 4}

$\theta\mid x \sim \text{Beta}(5,3)$. Estimate $E(\theta\mid x)$ and $P(\theta > 0.6\mid x)$ using Monte Carlo.

\subsection*{4(a)}
\[
E(\theta\mid x) = \frac{5}{5+3} = \frac{5}{8}
\]
Set \texttt{theta = rbeta(U, 5, 3)}. Then
\[
E(\theta\mid x) \approx \frac{1}{U}\sum \theta^{(u)} = \texttt{mean(theta)}
\]
\[
P(\theta > 0.6\mid x) \approx \frac{1}{U}\sum \mathbf{1}\{\theta^{(u)} > 0.6\} = \texttt{mean(theta > 0.6)}
\]

\subsection*{4(b)}

% ---- PASTE R CODE FOR 4(b) BELOW ----
\begin{verbatim}
# (b) U = 1000
set.seed(465)
U <- 1000
theta <- rbeta(U, 5, 3)
mean(theta)            # estimate of E(theta | x)
mean(theta > 0.6)      # estimate of P(theta > 0.6 | x)
\end{verbatim}

\[
\hat{E}(\theta\mid x) = 0.6232, \qquad \hat{P}(\theta > 0.6\mid x) = 0.585
\]

\subsection*{4(c)}

% ---- PASTE R CODE FOR 4(c) BELOW ----
\begin{verbatim}
# (c) Monte Carlo standard error and 95% MC confidence interval
mcse <- sd(theta) / sqrt(U)
sd(theta); mcse
mean(theta) + c(-1, 1) * 1.96 * mcse
\end{verbatim}

\[
\widehat{SE} = \frac{0.16483}{\sqrt{1000}} \approx 0.00521
\]
\[
0.6232 \pm 1.96(0.00521) = (0.6130,\ 0.6334)
\]

\subsection*{4(d)}

% ---- PASTE R CODE FOR 4(d) BELOW ----
\begin{verbatim}
# (d) U = 10000
set.seed(465)
U2 <- 10000
theta2 <- rbeta(U2, 5, 3)
mean(theta2)
mean(theta2 > 0.6)
mcse2 <- sd(theta2) / sqrt(U2)
sd(theta2); mcse2
mean(theta2) + c(-1, 1) * 1.96 * mcse2
mcse / mcse2           # ratio of standard errors (~ sqrt(10))

# exact values for comparison
5 / 8
1 - pbeta(0.6, 5, 3)
\end{verbatim}

The MC error dropped from $0.00521 \to 0.00162$, a factor of about 3.21. Since $SE = \frac{\sigma}{\sqrt{U}}$, multiplying $U$ by 10 would shrink the error by $\sqrt{10}$.

\subsection*{4(e)}
A 95\% posterior credible interval for $\theta$ describes the uncertainty about the parameter $\theta$ given the data. The width depends on the data and prior, so it doesn't shrink as $U$ grows.

A 95\% Monte Carlo CI for $E(\theta\mid x)$ describes simulation error only. The interval reflects how precisely our finite sample of draws estimates it. Its width scales like $\frac{1}{\sqrt{U}}$ and goes to 0 as $U \to \infty$.

%==================================================
\section*{Question 5}

$f(x) = 2x$, \quad $0 < x < 1$

\subsection*{5(a)}
$f(x) = 2x \geq 0$ on $(0,1)$.
\[
\int_0^1 2x\,dx = \left[x^2\right]_0^1 = (1)^2 - (0)^2 = 1 \;\Rightarrow\; \text{valid}
\]
\[
E_f(X) = \int_0^1 x\cdot 2x\,dx = \left[\frac{2x^3}{3}\right]_0^1 = \frac{2}{3}
\]

\subsection*{5(b)}
$g(x) = 1$ for $0 < x < 1$.
\[
2x \leq M,\ x\in(0,1) \;\Rightarrow\; \sup_{0<x<1}\frac{f(x)}{g(x)} = 2
\]
The constant $M = 2$ is valid and is the smallest valid envelope constant.

\subsection*{5(c)}
\begin{enumerate}
  \item Draw $X^* \sim \text{Unif}(0,1)$ and independent $U \sim \text{Unif}(0,1)$.
  \item Acceptance ratio: $\dfrac{f(X^*)}{Mg(X^*)} = \dfrac{2X^*}{2\cdot 1} = X^*$
  \item If $U \leq X^*$, accept and keep $X^*$ as a draw from $f$; otherwise reject and repeat.
\end{enumerate}
Theoretical acceptance probability:
\[
P(\text{accept}) = \int_0^1 \frac{f(x)}{Mg(x)}\,g(x)\,dx = \frac{1}{M} = \frac{1}{2}
\]
So about 2 proposals on average per accepted draw.

\subsection*{5(d)}

% ---- PASTE R CODE FOR 5(d) BELOW ----
\begin{verbatim}
# (d) Acceptance-rejection with g = Unif(0,1), M = 2
set.seed(465)
n_target <- 5000
draws <- numeric(n_target)
n_acc <- 0; n_try <- 0
while (n_acc < n_target) {
  x_star <- runif(1)
  u <- runif(1)
  n_try <- n_try + 1
  if (u <= x_star) {          # accept if u <= f(x*) / (M g(x*)) = x*
    n_acc <- n_acc + 1
    draws[n_acc] <- x_star
  }
}
mean(draws)                   # compare with exact E_f(X) = 2/3
n_acc / n_try                 # empirical acceptance rate (theory: 1/2)
\end{verbatim}

The sample mean of the 5000 accepted draws is 0.6647, close to the exact $E_f(X) = \frac{2}{3}$. Our difference of 0.002 is within about 1 Monte Carlo SE. It took 10,063 proposals, an acceptance rate of 0.497, which is close to our $\frac{1}{2}$.

\subsection*{5(e)}
This estimator doesn't need an acceptance step since it never tries to produce draws distributed exactly as $f$. Every draw from our $g$ is kept, and our mismatch between $g$ and $f$ is corrected by weighting each draw by
\[
w(z) = \frac{f(z)}{g(z)}
\]
Draws where $f$ is large count more and draws where $f$ is small count less. Acceptance-rejection instead fixes the mismatch by throwing away proposals, which here discards about half of them.

\end{document}